%=====================================================================
%  ch02 —— 原文 497–647：§2 群代数与卷积代数（2.1 概览、2.2 卷积与群代数）
%=====================================================================

\section{群代数与卷积代数的表示}
\label{s:repthryfingroup}
\subsection{概览}

\ssszh 设 $G$ 是整数环上的分裂约化群，\footnote{读者若通篇把 $G$ 取为可逆
$n \times n$ 矩阵之群 $GL_n$，甚至只取 $GL_2$，损失也不大。} 它在有限域、局部域与实数上
的点，其复表示论之间有富于启发的类比。借此可以把构造与想法从一个领域搬到另一个领域。
Harish-Chandra 甚至把这种彼此渗透称作表示论的 Lefschetz 原理
\cite{langlands1985harish}。\footnote{回想通常的 Lefschetz 原理：对任意复代数簇成立的
事实，例如一般光滑性，对任意特征零的代数闭域上的代数簇也成立。}

类似地，范畴理论中有一条基本原理，其源头是 Drinfeld 与 Laumon 关于整体几何朗兰兹对应的
工作，以及 Lusztig 关于特征标层的工作。这条原理说：$G(\mathbb{C})$ 的范畴表示论，其行为
应当与 $G(\mathbb{F}_q)$ 的复表示论（的一部分）相似；$G(\mathbb{C}(\!(t)\!))$ 的范畴
表示论，其行为应当与 $G(\mathbb{F}_q(\!(t)\!))$ 的复表示论相似。

\ssszh 有鉴于此，在进入范畴情形之前，我们先回顾通常有限群表示论的若干要点，然后把讨论
特殊化到 $G(\mathbb{F}_q)$。

\subsection{通过卷积看群代数}

\ssszh 设 $H$ 是有限群，$k$ 是特征零的代数闭域。用 $\on{Rep}(H)^\heartsuit$ 记 $H$ 在
$k$-向量空间上的表示所构成的阿贝尔范畴。

\begin{re}
$\on{Rep}(H)^\heartsuit$ 上的上标，是为了强调我们说的是阿贝尔范畴，即相应导出范畴
$\on{Rep}(H)$ 上标准 $t$-结构的核。类似地，下文中出现的模范畴与层范畴，除另有说明外
都指导出版本；相应的阿贝尔范畴则以上标「$\heartsuit$」标记。
\end{re}

\ssszh 回忆：给 $H$ 在向量空间 $V$ 上的一个表示，就是给一个映射
%
\begin{equation} \label{e:hom}
\alpha: H \rightarrow \on{End}(V).
\end{equation}
%
这里 $\on{End}(V)$ 指 $V$ 的 $k$-线性自同态，而 $\alpha$ 是单位的，并把 $H$ 中的乘法与
$\on{End}(V)$ 中的复合对应起来。更一般地，对任意 $k$-代数 $A$，考虑半群映射的集合
%
\begin{equation}  \label{e:homalg}
H \rightarrow A,
\end{equation}
%
其中 $A$ 以乘法看作半群。一个关键的观察是：由于我们丢掉了 $A$ 上的结构，即加法，
这可以扩充为从由 $H$ 造出的 $k$-代数（即群代数）出发的映射。下面我们换一种说法来讲群代数
的构造，以便适应后文要遇到的范畴变体。

\ssszh \label{ss:opsfinitesets}对有限集 $X$，用 $\on{Fun}(X)$ 记 $X$ 上 $k$-值函数
构成的向量空间。回忆这个对应有两条基本性质。第一，$X$ 分解为乘积 $X = X_1 \times X_2$，
给出其函数空间的一个典范张量积分解
$$ \Fun(X) \simeq \Fun(X_1) \otimes \Fun(X_2).$$
第二，给定有限集之间的映射 $f: X \rightarrow Y$，有拉回与推前映射
%
$$
 f^!: \on{Fun}(Y) \rightarrow \Fun(X) \quad \text{与} \quad
f_*: \Fun(X) \rightarrow \Fun(Y).
$$
%
具体地说，$f^!$ 把映射 $\phi: Y \rightarrow k$ 送到复合 $\phi \circ f$；而 $f_*$ 沿
$f$ 的纤维对 $\Fun(X)$ 中的函数 $\psi$ 求积分，即对 $Y$ 的任意元素 $y$，令
%
\[
(f_*\psi)(y) := \underset{x \in f^{-1}(y)}{\Sigma} \hspace{.5mm} \psi(x).
\]
%

\ssszh 把这两条性质合起来，$H$ 上的群结构就使 $\Fun(H)$ 具有 $k$-代数结构。确切地说，
若记
%
$
m: H \times H \rightarrow H
$
%
为乘法映射，便得到函数上的卷积映射
%
\begin{equation} \label{e:convfun}
\Fun(H) \otimes \Fun(H) \simeq \Fun(H \times H) \xrightarrow{ m_*}
\Fun(H).
\end{equation}
%

具体地说，若对 $H$ 的元素 $h$ 用 $\delta_h$ 记 $\Fun(H)$ 中相应的 delta 函数，即
%
\[
\delta_h(h') := \begin{cases}   1 & \text{若 } h = h', \\ 0 & \text{否则,}
\end{cases}
\]
%
则 $\delta_h$（$h \in H$）构成 $\Fun(H)$ 的一组基，并按规则
%
\[
\delta_{h_1} \star \delta_{h_2} = \delta_{h_1h_2}
\]
%
相乘，其中 $h_1h_2$ 指 $H$ 中的乘积。对应 $h \mapsto \delta_h$ 给出映射
$H \rightarrow \Fun(H)$，其性质是：任何形如 \eqref{e:homalg} 的映射，都唯一地扩充为
$k$-代数映射
%
\[
\Fun(H) \rightarrow A.
\]
%
特别地，这给出 $H$ 的表示与 $\Fun(H)$ 的左模之间熟悉的等价。

显然，$\Fun(H)$ 由 $H$ 的点的线性组合构成；在 $H$ 于向量空间上的表示中，它按相应的
自同态的线性组合作用。
